\documentclass[11pt,a4paper]{article}
\usepackage[bahasai]{babel}
\usepackage{amsmath, amssymb}
\usepackage{geometry}
\geometry{top=2.5cm, bottom=2.5cm, left=2.5cm, right=2.5cm}
\usepackage[most]{tcolorbox}
\usepackage{enumitem}
\usepackage{array}

% Semua kotak hitam-putih (hemat tinta)
\tcbset{colback=white, colframe=black, colbacktitle=white, coltitle=black,
  fonttitle=\bfseries, boxrule=0.6pt, arc=1mm,
  before upper={\setlength{\parskip}{5pt}}}

\setlength{\parindent}{0pt}
\setlength{\parskip}{5pt}

% Kotak "Sekarang Giliranmu" (garis putus-putus, ada ruang menulis)
\newtcolorbox{giliranmu}{enhanced, borderline={0.6pt}{0pt}{black, dashed},
  frame hidden, title={Sekarang Giliranmu!}, attach boxed title to top left={xshift=3mm, yshift=-2mm},
  boxed title style={colframe=white, colback=white}}

\newcommand{\kosong}{\rule{1.2cm}{0.4pt}}

% Kotak abu-abu (monokrom)
\newtcolorbox{selidiki}{enhanced, breakable, colback=black!4, colframe=black!60,
  colbacktitle=black!15, coltitle=black, title={Ayo Selidiki!}}
\newtcolorbox{contoh}{enhanced, breakable, colback=black!4, colframe=black!55,
  boxrule=0pt, leftrule=3pt, arc=0pt, left=3mm, before skip=6pt, after skip=6pt}
\newtcolorbox{rangkuman}{enhanced, colback=white, colframe=black!60,
  colbacktitle=black!15, coltitle=black, title={Rangkuman}}

% Label level: blok abu-abu kecil
\newcommand{\level}[1]{\par\medskip\noindent\colorbox{black!12}{\strut\textbf{#1}}\par\smallskip}

\begin{document}

\section*{Topik 2: Pola Angka Satuan}

% --------------------------------------------------------------------
\begin{selidiki}
Ambil pensilmu. Lengkapi tabel di bawah ini. Kamu cukup menulis \textbf{angka satuannya}
saja (angka paling belakang).

\begin{center}
\renewcommand{\arraystretch}{1.6}
\begin{tabular}{|l|*{8}{>{\centering\arraybackslash}p{0.9cm}|}}
\hline
Pangkat ke- & 1 & 2 & 3 & 4 & 5 & 6 & 7 & 8\\ \hline
Satuan $2^n$ & 2 & 4 & 8 & & & & & \\ \hline
Satuan $3^n$ & 3 & 9 & & & & & & \\ \hline
\end{tabular}
\end{center}

\textit{Petunjuk: kamu tidak perlu menghitung $2^8$ sampai habis. Cukup kalikan angka satuan
sebelumnya. Misalnya, satuan $2^3$ adalah 8, jadi satuan $2^4$ sama dengan satuan
$8 \times 2 = 16$, yaitu 6.}

Sekarang jawab:
\begin{enumerate}
    \item Apa yang kamu lihat pada baris $2^n$? Angka apa saja yang muncul berulang?
    \item Setelah berapa langkah polanya kembali ke awal?
    \item Tanpa menghitung, tebak angka satuan $2^{9}$ dan $2^{10}$.
\end{enumerate}

Kalau kamu menemukan bahwa angka satuannya berulang terus, selamat! Kamu baru saja menemukan
rahasia untuk menghitung bilangan raksasa seperti $3^{2026}$ tanpa kalkulator.
\end{selidiki}

% --------------------------------------------------------------------
\subsection*{Langkah 1: Bagi pangkatnya dengan panjang pola}
Pola angka satuan dari pangkat 3 adalah \textbf{3, 9, 7, 1}, lalu berulang. Panjang polanya 4.

\begin{contoh}
\textbf{Contoh.} Berapa angka satuan dari $3^{2026}$?

\textit{Jawab:} Bagi pangkatnya dengan 4: $2026 \div 4 = 506$ sisa \textbf{2}. Ambil angka
urutan ke-2 dari pola (3, \textbf{9}, 7, 1). Jadi angka satuannya \textbf{9}.
\end{contoh}

\begin{giliranmu}
\begin{enumerate}[label=\alph*.]
    \item Angka satuan dari $2^{15}$ adalah \kosong. \quad (pola 2: 2, 4, 8, 6)
    \item Angka satuan dari $7^{42}$ adalah \kosong. \quad (pola 7: 7, 9, 3, 1)
\end{enumerate}
\end{giliranmu}

% --------------------------------------------------------------------
\subsection*{Langkah 2: Hati-hati kalau sisanya 0}
Kalau pangkatnya habis dibagi panjang pola, tidak ada ``urutan ke-0''. Ambil angka yang
\textbf{paling terakhir} dari pola.

\begin{contoh}
\textbf{Contoh.} Berapa angka satuan dari $3^{100}$?

\textit{Jawab:} $100 \div 4 = 25$ sisa \textbf{0}. Ambil angka terakhir dari pola
(3, 9, 7, \textbf{1}). Jadi angka satuannya \textbf{1}.
\end{contoh}

\begin{giliranmu}
\begin{enumerate}[label=\alph*.]
    \item Angka satuan dari $8^{100}$ adalah \kosong. \quad (pola 8: 8, 4, 2, 6)
    \item Angka satuan dari $2^{2024}$ adalah \kosong.
\end{enumerate}
\end{giliranmu}

% --------------------------------------------------------------------
\subsection*{Langkah 3: Angka di depan tidak penting}
Angka satuan hasil perkalian hanya ditentukan oleh angka satuan bilangannya. Jadi angka
ribuan, ratusan, dan puluhan boleh kita abaikan.

\begin{contoh}
\textbf{Contoh.} Berapa angka satuan dari $2027^{2026}$?

\textit{Jawab:} Angka satuan 2027 adalah 7, jadi cukup cari angka satuan $7^{2026}$.
$2026 \div 4$ sisa 2. Urutan ke-2 dari pola (7, \textbf{9}, 3, 1) adalah \textbf{9}.
\end{contoh}

\begin{giliranmu}
\begin{enumerate}[label=\alph*.]
    \item Angka satuan dari $1234^{5}$ adalah \kosong.
    \item Angka satuan dari $2029^{2026}$ adalah \kosong.
\end{enumerate}
\end{giliranmu}

% --------------------------------------------------------------------
\subsection*{Langkah 4: Menjumlah dan mengali}
Cari angka satuan masing-masing dulu, baru dijumlah atau dikali. Setelah itu, ambil lagi
angka satuannya.

\begin{contoh}
\textbf{Contoh.} Berapa angka satuan dari $4^{2025} + 6^{2025}$?

\textit{Jawab:} Pola 4 adalah 4, 6, 4, 6, \dots{} Pangkatnya ganjil, jadi satuannya 4.
Pangkat dari 6 selalu berakhiran 6. Jadi $4 + 6 = 10$, dan angka satuannya \textbf{0}.
\end{contoh}

\begin{giliranmu}
\begin{enumerate}[label=\alph*.]
    \item Angka satuan dari $2^{2025} + 2^{2026}$ adalah \kosong.
    \item Angka satuan dari $2^{10} \times 3^{10}$ adalah \kosong.
\end{enumerate}
\end{giliranmu}

% --------------------------------------------------------------------
\begin{rangkuman}
\begin{enumerate}
    \item Tulis beberapa angka satuan awal sampai polanya berulang.
    \item Bagi pangkatnya dengan panjang pola, lalu lihat sisanya.
    \item Sisa 1 berarti urutan ke-1, sisa 2 berarti urutan ke-2, dan seterusnya.
          \textbf{Sisa 0 berarti urutan terakhir.}
    \item Angka di depan tidak penting. Cukup lihat angka satuannya.
    \item Untuk menjumlah atau mengali, cari angka satuan masing-masing dulu.
\end{enumerate}
\end{rangkuman}

\textit{Kunci Giliranmu:} Langkah 1: 8, 9. \quad Langkah 2: 6, 6. \quad Langkah 3: 4, 1. \quad
Langkah 4: 6, 6.

\newpage
% --------------------------------------------------------------------
\subsection*{Latihan Soal}

\level{Level 1: Pemanasan}
\begin{enumerate}
    \item Berapakah angka satuan dari $9^{2026}$?
    \item \textbf{(Bakteri Super)} Di sebuah laboratorium, sebuah sel bakteri membelah diri
          menjadi 3 sel baru setiap jamnya. Pada jam pertama ada 3 bakteri, jam kedua ada
          9 bakteri, dan seterusnya. Berapakah angka terakhir dari total bakteri pada jam ke-50?
    \item \textbf{(Mesin Tiket)} Sebuah mesin pencetak tiket otomatis mengeluarkan nomor seri
          dengan aturan $7^n$, dengan $n$ adalah urutan pembeli. Pembeli ke-41 akan mendapatkan
          tiket dengan nomor seri yang berakhiran angka berapa?
    \item Berapakah angka satuan dari $5^{2026}$ dan $6^{2026}$?
    \item Berapakah angka satuan dari $4^{99}$?
\end{enumerate}

\level{Level 2: Eksplorasi}
\begin{enumerate}
    \setcounter{enumi}{5}
    \item Tentukan angka satuan dari hasil $3^{10} + 7^{10}$.
    \item \textbf{(Gembok Koper Budi)} Budi lupa kombinasi 3 digit gembok kopernya. Ia hanya
          ingat sebuah petunjuk yang ia tulis di bukunya: \textit{``Digit pertama adalah angka
          satuan dari $2^{40}$, digit kedua angka satuan dari $3^{41}$, dan digit ketiga angka
          satuan dari $4^{42}$''}. Berapakah kombinasi koper Budi?
    \item \textbf{(Sandi Rahasia Agen 007)} Seorang agen rahasia harus memasukkan pin 1 digit
          untuk membuka pintu brankas. Petunjuk pin tersebut berbunyi: \textit{``Angka satuan dari
          jumlah hari dalam seminggu dipangkatkan dengan tahun kemerdekaan Indonesia''}.
          Angka berapakah yang harus ditekan agen tersebut?
    \item Tentukan angka satuan dari hasil pengurangan $9^{99} - 4^{44}$.
    \item \textbf{(Kalkulator Rusak)} Tombol layar kalkulator Andi rusak sehingga hanya bisa
          menampilkan 1 digit paling belakang dari sebuah perhitungan. Jika Andi mengetikkan
          $2^{2025} + 2^{2026} + 2^{2027}$, angka berapakah yang muncul di layar?
\end{enumerate}

\level{Level 3: Tantangan OSK}
\begin{enumerate}
    \setcounter{enumi}{10}
    \item Tentukan angka satuan dari $2023^{2024} + 2024^{2023}$.
    \item \textbf{(Harta Karun Raja)} Seorang Raja memberikan hadiah koin emas kepada ksatria
          terbaiknya. Jumlah koin yang diberikan adalah sebanyak hasil perkalian
          $13^{13} \times 17^{17}$. Ksatria tersebut ingin membagikan koinnya kepada 10 panti
          asuhan sama rata. Berapakah koin yang tersisa (tidak habis terbagi)?
    \item Tentukan angka satuan dari $1^1 + 2^2 + 3^3 + 4^4 + 5^5$.
    \item \textbf{(Bom Waktu)} Sebuah mesin pemecah sandi harus menebak digit terakhir dari
          $2^{3^4}$. Angka berapa yang harus dimasukkan? (Kerjakan pangkat yang ada di atas dulu ya!)
    \item Tentukan angka satuan dari $1^{2026} + 2^{2026} + 3^{2026} + \dots + 9^{2026}$.
\end{enumerate}

\level{Level 4: Setara OSP (Provinsi), tulis jawaban lengkap dengan alasannya}
\begin{enumerate}
    \setcounter{enumi}{15}
    \item Buktikan bahwa untuk setiap bilangan asli $n$, bilangan $n^5 - n$ selalu habis
          dibagi 10.
\end{enumerate}

\level{Level 5: Setara OSN (Nasional), tulis jawaban lengkap dengan alasannya}
\begin{enumerate}
    \setcounter{enumi}{16}
    \item Tentukan semua bilangan asli $n$ sehingga $n^n$ berakhiran angka 7.
          Buktikan bahwa tidak ada bilangan lain yang memenuhi.
\end{enumerate}

\end{document}
